\section{Modello Asmeta}

\subsection{Il modello (\texttt{asmeta/src/SmartParking.asm})}

Nel modello ho definito tre domini enumerativi (\texttt{StatoSistema},
\texttt{TipoUtente}, \texttt{TipoPosto}), sei funzioni monitorate (\texttt{sens\_in},
\texttt{sens\_out}, \texttt{transito\_ok}, \texttt{auto\_via}, \texttt{utente\_rilevato},
\texttt{pagamento\_ok}) e quattro funzioni controllate (\texttt{stato},
\texttt{posti\_std}, \texttt{posti\_dis}, \texttt{posto\_assegnato}). A queste ho
aggiunto una funzione derivata n-aria (binaria), \texttt{assegnabile: Prod(TipoUtente,
TipoPosto) $\to$ Boolean}, che raccoglie in un punto solo tutta la politica di
assegnazione: un \texttt{POSTO\_STD} si può dare a chiunque se \texttt{posti\_std > 0},
mentre un \texttt{POSTO\_DIS} è riservato ai disabili e solo se \texttt{posti\_dis > 0}.
Trattandosi di una funzione derivata, viene ricalcolata a ogni stato a partire dalle
funzioni controllate, quindi non occupa memoria.

\begin{lstlisting}[language=Pascal, caption={Funzione derivata assegnabile}]
derived assegnabile: Prod(TipoUtente, TipoPosto) -> Boolean

function assegnabile($u in TipoUtente, $p in TipoPosto) =
    ($p = POSTO_STD and posti_std > 0) or
    ($p = POSTO_DIS and $u = DISABILE and posti_dis > 0)
\end{lstlisting}

Sempre tra le derivate ho aggiunto \texttt{parcheggioPieno}, che uso per far vedere il
costrutto \texttt{forall}: è vera quando, per \emph{ogni} tipo di utente, non è
assegnabile né un posto standard né uno disabile, cioè quando davvero non può più entrare
nessuno. La controllo nello scenario di saturazione \texttt{test\_pieno\_totale.avalla},
dove all'inizio è falsa (c'è ancora posto) e a parcheggio pieno diventa vera.

\begin{lstlisting}[language=Pascal, caption={Funzione derivata parcheggioPieno (uso di forall)}]
function parcheggioPieno =
    (forall $u in TipoUtente with
        (not assegnabile($u, POSTO_STD) and not assegnabile($u, POSTO_DIS)))
\end{lstlisting}

La regola più importante è \texttt{r\_gestione\_VER}, quella che decide a chi assegnare
il posto. Qui ho usato una regola \texttt{let} insieme alla derivata
\texttt{assegnabile}: la regola prova prima con \texttt{POSTO\_DIS} (che è vero solo per
un disabile quando ci sono ancora posti disabili liberi), poi con \texttt{POSTO\_STD}
(il caso di STD e abbonati, ma anche il fallback del disabile quando i posti riservati
sono finiti) e, se non va a buon fine nessuno dei due, nega l'accesso mandando il
sistema in NEG:

\begin{lstlisting}[language=Pascal, caption={Regola r\_gestione\_VER (assegnazione posto, con let e derivata)}]
rule r_gestione_VER =
    let ($u = utente_rilevato) in
        if assegnabile($u, POSTO_DIS) then
            par
                stato := INGR
                posto_assegnato := POSTO_DIS
            endpar
        else
            if assegnabile($u, POSTO_STD) then
                par
                    stato := INGR
                    posto_assegnato := POSTO_STD
                endpar
            else
                stato := NEG
            endif
        endif
    endlet
\end{lstlisting}

Per il model checking ho poi usato una versione semplificata del modello
(\texttt{SmartParking\_MC.asm}). Qui ho lasciato gli if annidati "appiattiti", che si
comportano esattamente allo stesso modo ma rendono più facile la traduzione verso
NuSMV. Sempre per lo stesso motivo, in questo modello i contatori non sono più
\texttt{Integer} ma stanno su un dominio finito \texttt{Posti = \{0..2\}}, dato che
AsmetaSMV non accetta domini infiniti. Ho scelto \texttt{\{0..2\}} e non \texttt{\{0,1\}}
apposta: così le proprietà \texttt{ag(posti\_std <= 1)} e \texttt{ag(posti\_dis <= 1)}
sono verifiche vere e proprie, e non diventano ovvie solo per come è fatto il tipo.

\subsubsection{Simulazione manuale}

Prima di lanciare gli scenari automatici, per prendere un po' di confidenza con il
modello, conviene simulare a mano qualche step con il simulatore di Asmeta.

\screenshot{asmeta_simulator_ingresso_std}{Simulatore -- ingresso STD completo}{%
Aprire \texttt{asmeta/src/SmartParking.asm} in Eclipse e lanciare il simulatore
interattivo. Impostare \texttt{sens\_in=true}, fare Step, poi \texttt{utente\_rilevato=STD}
e proseguire (CHK\_IN$\to$VER, VER$\to$INGR), infine \texttt{transito\_ok=true} e Step.
Lo screenshot deve mostrare la sequenza di stati IDLE$\to$CHK\_IN$\to$VER$\to$INGR$\to$IDLE
con \texttt{posti\_std} che passa da 1 a 0. La colonna State~0 documenta anche lo stato
iniziale (\texttt{stato=IDLE}, sensori a false); si noti che l'animatore mostra le
funzioni solo dal primo stato in cui vengono lette o aggiornate.}

\screenshot{asmeta_simulator_uscita_std}{Simulatore -- uscita STD con pagamento}{%
Proseguendo nella stessa sessione: \texttt{sens\_out=true} $\to$ CHK\_OUT, utente STD
$\to$ TARIF, \texttt{pagamento\_ok=true} $\to$ USC, \texttt{transito\_ok=true} $\to$
IDLE. Lo screenshot deve mostrare il passaggio per TARIF (pagamento) e lo stato finale
con \texttt{posti\_std} tornato a 1 e \texttt{posto\_assegnato = NESSUNO}.}

\subsection{Scenari Avalla}

\begin{longtable}{@{}p{4.2cm}p{8.8cm}@{}}
\toprule
\textbf{Scenario} & \textbf{Obiettivo} \\
\midrule
\endhead
\texttt{test\_abbonato.avalla} & Ciclo completo ABBONATO: ingresso senza pagamento, uscita diretta (salta TARIF). Bug originale corretto: init \texttt{posti\_std} non era 5 ma 1, e lo stato \texttt{PARK} non esiste (è \texttt{IDLE}). \\
\texttt{test\_disabile.avalla} & Ciclo DISABILE: assegnazione POSTO\_DIS, \texttt{posto\_assegnato} resta POSTO\_DIS durante la sosta, azzerato solo in uscita. \\
\texttt{test\_pieno.avalla} & Un utente STD occupa l'unico posto; un secondo STD viene respinto (NEG). \\
\texttt{test\_disabile\_fallback.avalla} & Copre il fallback: un secondo DISABILE, con \texttt{posti\_dis=0}, riceve POSTO\_STD. \\
\texttt{test\_pagamento\_std.avalla} & Ciclo di uscita STD con un pagamento fallito (resta in TARIF) seguito da uno riuscito. \\
\texttt{test\_pieno\_totale.avalla} & Parcheggio completamente saturo (STD + DISABILE), un terzo utente DISABILE viene respinto. \\
\texttt{test\_random\_30steps.avalla} & Scenario \emph{generato automaticamente}: 30 step random eseguiti nel simulatore ed esportati con ``export to avalla'' (saturazione del parcheggio, doppio NEG, ciclo di pagamento TARIF, rientro). \\
\bottomrule
\end{longtable}

\screenshot{asmeta_animator_test_abbonato}{Animazione di test\_abbonato.avalla}{%
Eseguire lo scenario nell'animatore Asmeta. La tabella degli stati deve mostrare la
fase di uscita dell'ABBONATO: CHK\_OUT $\to$ USC \emph{senza passare da TARIF}
(nessun pagamento), con \texttt{posti\_std} che torna a 1 e \texttt{posto\_assegnato}
azzerato a NESSUNO. L'esito dei \texttt{check} (tutti PASSED) è riportato dalla
validazione AsmetaV in console.}

\screenshot{asmeta_animator_test_disabile_fallback}{Animazione di test\_disabile\_fallback.avalla}{%
Stessa procedura per lo scenario di fallback. Deve mostrare i due ingressi DISABILE
consecutivi con l'assegnazione POSTO\_DIS poi POSTO\_STD (fallback), tutti i check verdi.}

\screenshot{asmeta_animator_test_pieno_totale}{Animazione di test\_pieno\_totale.avalla}{%
Stessa procedura per lo scenario di saturazione totale, con il rifiuto (NEG) del
terzo utente e tutti i check verdi.}

\subsection{AsmetaV -- validazione con copertura (Vc)}\label{sec:asmetav-vc}

Gli scenari li ho validati con AsmetaV usando la modalità che calcola anche la
copertura, cioè l'opzione ``Vc'' richiesta dalla consegna. In pratica, per ogni scenario
il validatore non si limita a controllare i \texttt{check}, ma stampa in console anche
quali regole del modello sono state usate (la \emph{rule coverage}). Mettendo insieme i
sei scenari scritti a mano si arriva a coprire tutte le 8 regole di transizione, la main
rule e tutti i rami di \texttt{r\_gestione\_VER}.

A questo punto è interessante confrontare questi scenari con quelli generati a caso. Un
primo export di 15 step random copriva solo 6 regole su 9 (restavano fuori
\texttt{r\_gestione\_NEG}, \texttt{r\_gestione\_TARIF} e \texttt{r\_gestione\_USC});
raddoppiando a 30 step (\texttt{test\_random\_30steps.avalla}) la rule coverage arriva
al 100\%, ma la branch coverage resta comunque incompleta
(\texttt{r\_gestione\_USC} 33\%, \texttt{r\_gestione\_TARIF} 50\%,
\texttt{r\_gestione\_IDLE} 75\%). Il ramo ``pagamento fallito'' di TARIF, per esempio,
non capita quasi mai per caso, e infatti viene coperto solo dallo scenario mirato
\texttt{test\_pagamento\_std}. Insomma, la generazione random è utile come aggiunta, ma
non sostituisce gli scenari scritti a mano sui casi limite.

\screenshot{asmetav_coverage_report}{Report di copertura AsmetaV (Vc)}{%
Tasto destro su ciascun file \texttt{.avalla} $\to$ \textbf{ASMETA $\to$ Validate
scenario and compute coverage} (voce con ``coverage''). Lo screenshot deve mostrare la
console con l'esito dei check (PASSED) e l'elenco delle regole coperte; ripetere per
tutti gli scenari e mostrare (anche cumulativamente) che tutte le regole risultano
coperte.}

\subsection{Model checking (\texttt{asmeta/src/SmartParking\_MC.asm})}

Per il model checking ho scritto 10 CTLSPEC (Tabella~\ref{tab:ctlspec}). Otto me le
aspetto vere e le ho divise in proprietà di sicurezza (safety), di raggiungibilità e di
liveness; le altre due invece le ho messe apposta false, giusto per far vedere che il
model checker, quando una proprietà non regge, tira fuori un controesempio.

\begin{longtable}{@{}p{0.6cm}p{7.3cm}p{2.2cm}p{1.3cm}@{}}
\toprule
\textbf{\#} & \textbf{CTLSPEC} & \textbf{Categoria} & \textbf{Esito} \\
\midrule
\endhead
1  & \texttt{ag(posti\_std >= 0)} & Safety & TRUE \\
2  & \texttt{ag(posti\_dis >= 0)} & Safety & TRUE \\
3  & \texttt{ag((stato=INGR and posto\_assegnato=POSTO\_STD) implies posti\_std>0)} & Safety & TRUE \\
4  & \texttt{ag(posti\_std <= 1)} & Safety & TRUE \\
5  & \texttt{ag(posti\_dis <= 1)} & Safety & TRUE \\
6  & \texttt{ag((stato=INGR and posto\_assegnato=POSTO\_DIS) implies posti\_dis>0)} & Safety & TRUE \\
7  & \texttt{ef(stato = TARIF)} & Raggiungibilità & TRUE \\
8  & \texttt{ag(stato=TARIF implies ef(stato=IDLE))} & Liveness & TRUE \\
9  & \texttt{ag(stato != NEG)} & Falsa (voluta) & FALSE \\
10 & \texttt{ag(stato=TARIF implies af(stato=IDLE))} & Falsa (voluta) & FALSE \\
\bottomrule
\caption{Le 10 CTLSPEC sottoposte ad AsmetaSMV: 8 verificate, 2 volutamente false.}
\label{tab:ctlspec}
\end{longtable}

Le due specifiche false sono interessanti, ognuna a modo suo:

\begin{itemize}
    \item la \#9 dice che ``l'accesso non viene mai negato''. Il controesempio che
    genera NuSMV (9 stati) fa entrare un abbonato che occupa l'unico posto standard;
    poi arriva un secondo veicolo, la verifica non trova posto e il sistema va in NEG,
    e lo fa con il posto disabili ancora libero (\texttt{posti\_dis = 1}). Questa
    traccia dimostra due cose in una volta: che lo stato di rifiuto si può davvero
    raggiungere, e che la riserva funziona, perché il posto disabili non viene dato a
    chi non ne ha diritto;
    \item la \#10 è la versione con \texttt{af} della liveness \#8: chiede che da TARIF
    il sistema torni in IDLE su \emph{ogni} cammino, non solo che esista un cammino che
    ci riporta. È falsa perché nulla vieta all'ambiente di tenere
    \texttt{pagamento\_ok = false} per sempre, cioè un utente che non paga mai, e infatti
    il controesempio è proprio il loop infinito in TARIF. Mettere a confronto la \#8 e la
    \#10 fa vedere bene la differenza tra \texttt{ef} e \texttt{af} in CTL.
\end{itemize}

\screenshot{asmetasmv_ctlspec_results}{Esito delle 10 CTLSPEC in AsmetaSMV}{%
Tasto destro su \texttt{SmartParking\_MC.asm} $\to$ \textbf{Run As $\to$ AsmetaSMV}.
Lo screenshot deve mostrare le 8 specifiche con esito \texttt{is true} e le 2
specifiche volutamente false con esito \texttt{is false}, ciascuna con il proprio
controesempio generato.}

\screenshot{asmetasmv_random_trace}{Esecuzione random-step di verifica manuale}{%
Nella stessa vista AsmetaSMV, usare la funzione di simulazione random (random step)
per generare una traccia di 8-10 stati. Utile per mostrare a voce, all'orale, che il
modello si comporta coerentemente (nessuna transizione viola le regole).}

\subsubsection{Model Advisor (opzionale)}

Il Model Advisor (AsmetaMA) lavora passando per la traduzione in NuSMV, quindi l'ho
lanciato sul modello semplificato \texttt{SmartParking\_MC.asm}: quello principale, con
\texttt{Integer} e funzioni derivate, non è traducibile e dà ``Error Domain Integer not
supported''. L'analisi delle meta-proprietà ha dato questi esiti:

\begin{itemize}
\item \textbf{MP1, MP3, MP4, MP7} non hanno prodotto segnalazioni: nessun aggiornamento
inconsistente, nessuna assegnazione banale, nessuna locazione rimuovibile o mai
aggiornata.
\item \textbf{MP2} (``conditional rule not complete'') scatta su tutte le regole di
transizione, ma si tratta di falsi positivi noti. L'advisor segnala ogni
\texttt{if} privo di \texttt{else}, mentre in ASM l'assenza dell'else è idiomatica e
significa ``nessun aggiornamento'': è proprio il comportamento che serve per i
self-loop della FSM (per esempio \texttt{r\_gestione\_NEG} resta in NEG finché
\texttt{auto\_via} non diventa vero, e \texttt{r\_gestione\_TARIF} resta in TARIF
a pagamento fallito).
\item \textbf{MP5/MP6} (``Domain Posti could be reduced: element \{2\} could be
removed'') non segnala un difetto ma conferma una scelta. Il dominio \texttt{Posti =
\{0..2\}} è stato preso apposta più ampio del necessario, e l'advisor
dimostra che il valore 2 non è mai raggiunto da \texttt{posti\_std}/\texttt{posti\_dis},
cioè la proprietà di capacità \texttt{ag(posti <= 1)} verificata anche dalle
CTLSPEC.
\end{itemize}

\screenshot{model_advisor_output}{Output del Model Advisor (estratto)}{%
Tasto destro su \texttt{SmartParking\_MC.asm} $\to$ \textbf{Model Advisor}. È
sufficiente un estratto rappresentativo della console: l'ideale è MP1 (nessun update
inconsistente), le prime righe di MP2 e le segnalazioni MP5/MP6 sul dominio Posti;
l'elenco completo delle MP2 è riassunto e discusso nel testo.}

\subsubsection{ATGT -- generazione automatica di scenari (opzionale)}

\screenshot{atgt_scenario_list}{Elenco degli scenari generati da ATGT}{%
Tasto destro su \texttt{SmartParking.asm} $\to$ \textbf{Run As $\to$ ATGT}. Mostrare
la lista dei file \texttt{.avalla} generati automaticamente nella cartella di output
(uno screenshot dell'Explorer/Package Explorer di Eclipse con la cartella espansa).}

Gli scenari che ATGT genera (\texttt{asmeta/src/abstractests/testtest*.avalla}) si
possono guardare direttamente nel repository. Il loro ruolo, insieme a quello dello
scenario random da 30 step, l'ho già discusso nel confronto con gli scenari scritti a
mano in \S\ref{sec:asmetav-vc}.
